\documentclass[main.tex]{subfiles}
\begin{document}
\section{Incomplete Information}\label{sec:incomplete_info}
So far, I have examined commitment games in which both players commit simultaneously. I now relax this assumption by introducing a game of incomplete information, where there is uncertainty over the order of commitment. 
%This approach not only captures situations in which players are unsure who commits first, but also those in which there is uncertainty about the recognition of commitment. Even if player $i$ commits before player $-i$, if $-i$ has not observed or inferred that $i$ has committed by the time she acts the strategic situation remains effectively simultaneous. The model thus also captures uncertainty about how quickly the other player will realize that a commitment has occurred.
At the outset, Nature draws one of three possible states:
\begin{itemize}
    \item $\omega_i$, player $i$ commits first, with probability $\mu$;
    \item $\omega_{-i}$, player $-i$ commits first, with probability $\mu$;
    \item $\omega_\emptyset$, both players commit simultaneously, with probability $1-2\mu$,
\end{itemize}
where $\mu \leq 0.5$. Note that when a player commits first in states $\omega_i$ and $\omega_{-i}$, the second-moving player still has the opportunity to commit before the base game is played, should the first mover choose to wait.
I model uncertainty as follows. When the state is $\omega_{-i}$, player $i$ moves secondand observes the state perfectly. This is because a player moving second must observe the other's prior commitment in order to best respond, and thus cannot face any uncertainty.
However, when player $i$ is called to act and observes no prior commitment, she cannot distinguish whether the true state is $\omega_i$ or $\omega_\emptyset$. 
Accordingly, I model the signal structure as follows: 
\[
\sigma_i(\omega_i)=\sigma_i(\omega_\emptyset)=t_\emptyset \quad\text{and}\quad \sigma_i(\omega_{-i})=t_{-i}.
\]
It follows that player $i$'s beliefs conditional on observing $t_\emptyset$ are given by $P_i(\omega_i|t_\emptyset)=\frac{\mu}{1-\mu}$ and $P_i(\omega_\emptyset|t_\emptyset)=\frac{1-2\mu }{1-\mu}$, while $P_i(\omega_{-i}|t_{-i})=1$. 
\subsection*{\centering  Equilibria}
I will now analyze the equilibria under this information structure using the stylized Game of Chicken from the introduction as the base game (see Table~\ref{tab:gameofchicken}).
 
When player $i$ receives the signal $t_{-i}$ (indicating player $i$ moves second), her equilibrium strategy is straightforward. She best responds to any commitment observed from the first mover. However, if she observes the first mover choosing to wait instead of committing, she will commit to going straight ($S$) to ensure her preferred equilibrium outcome. Now consider a player $i$ who receives the signal $t_\emptyset$. In this case, player $i$ has three possible strategies: commit to going straight ($S_1$), commit to moving out of the way ($M_1$), or wait ($W$). Assume player $-i$, when also receiving $t_\emptyset$, randomizes among these strategies with probabilities $p$, $q$, and $1 - p - q$, respectively. Then the payoffs for player $i$’s pure strategies are:
\begin{align*}
U_i(S_{1} \mid t_\emptyset)
  &= P_i(\omega_i \mid t_\emptyset) \cdot 1
     + P_i(\omega_\emptyset \mid t_\emptyset) \cdot
       \Bigl[\, p \cdot (-1) + q \cdot 1 + (1 - p - q) \cdot 1 \,\Bigr]
\\
U_i(M_{1} \mid t_\emptyset)
  &= P_i(\omega_i \mid t_\emptyset) \cdot \left(-\tfrac{1}{2}\right)
     + P_i(\omega_\emptyset \mid t_\emptyset) \cdot
       \Bigl[\, p \cdot \left(-\tfrac{1}{2}\right) + q \cdot 0 + (1 - p - q) \cdot \left(-\tfrac{1}{2}\right) \,\Bigr]
\\
U_i(W \mid t_\emptyset)
  &= P_i(\omega_i \mid t_\emptyset) \cdot \left(-\tfrac{1}{2}\right)
     + P_i(\omega_\emptyset \mid t_\emptyset) \cdot
       \Bigl[\, p \cdot \left(-\tfrac{1}{2}\right) + q \cdot 1 + (1 - p - q) \cdot \left(-\tfrac{1}{3}\right) \,\Bigr].
\end{align*}
Recall that $S_1$ denotes the strategy of committing to Straight upon receiving $t_\emptyset$, while $M_1$ denotes the strategy of committing to Move. Note that $M_1$ is weakly dominated by $W$.

\paragraph{Symmetric Pure Strategy Equilibrium} There is an equilibrium where both players commit to going straight after receiving $t_\emptyset$. Note that this corresponds to setting $p=1$ in the pure strategy payoffs above. Then:
\begin{equation}
    U_i(S_1 \mid t_\emptyset) = \frac{3\mu - 1}{1 - \mu}
    \quad \text{and} \quad
    U_i(W \mid t_\emptyset) = -\frac{1}{2},
    \label{eq:alwaysstraightaftert_emptyset}
\end{equation}
and $U_i(S_1 \mid t_\emptyset)\geq U_i(W \mid t_\emptyset)\Leftrightarrow\mu\geq0.2$. Thus, the symmetric subgame perfect Nash equilibrium with both players committing straight upon receiving $t_\emptyset$ exists when $\mu\geq0.2$. 
The expected equilibrium payoff is $U=P(t_\emptyset)\cdot\frac{3\mu - 1}{1 - \mu}+P(t_{-i})\cdot(-\frac{1}{2})$, which equals $2.5\mu-1$ since $P(t_\emptyset)=1-\mu$ and $P(t_{-i})=\mu$.
 
It is easy to see that there is no equilibrium in which both players wait after receiving $t_\emptyset$. A strategy of always waiting corresponds to $1-p-q=1$, but in that case $U_i(S_1 \mid t_\emptyset)> U_i(W \mid t_\emptyset)$, so waiting is not a best response to waiting.
\paragraph{Asymmetric Pure Strategies Equilibria}
There are also equilibria where one player commits to going straight after receiving $t_\emptyset$ and the other one waits. Let player~$-i$ be the one who commits straight. Then player~$i$'s payoffs are as in Equation~\ref{eq:alwaysstraightaftert_emptyset}, and waiting is a best response when $\mu \leq 0.2$. From player~$-i$'s perspective, $U_{-i}(S_1 \mid t_\emptyset) > U_{-i}(W \mid t_\emptyset)$ always holds when the other player waits ($1-p-q=1$), so she indeed prefers to commit straight. Hence, this is an equilibrium. Moreover, instead of waiting, player~$i$ could also choose to commit to moving out of the way ($M_1$) in equilibrium. These asymmetric pure strategy equilibria yield a payoff of $1$ to the player committing straight, and $-\tfrac{1}{2}$ to the player who waits (or moves aside).
\paragraph{Symmetric Mixed Equilibrium}
There is also a symmetric mixed equilibrium where players mix over $S_1$ and $W$ when receiving $t_\emptyset$. $M_1$ is weakly dominated by waiting and will not feature in a mixed equilibrium, thus $q=0$. The $p$ that equalizes $U_i(S_{1} \mid t_\emptyset)$ and $ U_i(W \mid t_\emptyset)$ is $p=\frac{8-7\mu}{11(1-2\mu)}$, which lies in $[0,1]$ for $\mu\leq0.2$. The expected equilibrium payoff is 
\begin{equation}
    U = P(t_\emptyset) \cdot U_i(S_1 \mid t_\emptyset)
      + P(t_{-i}) \cdot \left[\, p \cdot \left(-\tfrac{1}{2}\right) + (1 - p) \cdot 1 \,\right],
    \label{eq:expectedpayoff_incompleteinfo}
\end{equation}
which can be rearranged to $\frac{3\mu-5}{11}+\frac{-2\mu-23\mu^2}{22(1-2\mu)}$. Let us briefly zoom in on the dynamics in Equation~\ref{eq:expectedpayoff_incompleteinfo}. When the state is $\omega_i$ or $\omega_\emptyset$, player $i$ receives $t_\emptyset$ and mixes as described above. She then receives the expected payoff $U_i(S_1 \mid t_\emptyset)$ or $U_i(W \mid t_\emptyset)$, both are equal by the definition of mixed equilibria. 
When the state is $\omega_{-i}$, only $-i$ receives $t_\emptyset$ and mixes, while $i$ receives $t_{-i}$, knows that she can commit only as the second mover and observes the outcome of $-i$'s mixing. With probability $p$, that outcome is $S_1$, and $i$ as the second mover will move aside. However, with probability $1 - p$, the first mover waits, in which case $i$ will commit to going straight—inducing her preferred outcome.
\paragraph{Summary and Interpretation}
All symmetric equilibria are summarized in Table~\ref{tab:equilibrium-strategies}. 
\begin{table}[h!]
\centering
\begin{tabular}{@{}lll@{}}
\toprule
Condition & Strategy Profile & Expected Payoff \\
\midrule
$ \mu \leq 0.2 $ &
\begin{tabular}[t]{@{}l@{}}
$\forall i: \  \sigma_i(S_1 \mid t_\emptyset) = p,\ \sigma_i(W \mid t_\emptyset) = 1 - p,$ \\[0.5em]
where $ p = \tfrac{8 - 7\mu}{11(1 - 2\mu)} $
\end{tabular} &
$ \tfrac{3\mu - 5}{11} + \tfrac{-2\mu - 23\mu^2}{22(1 - 2\mu)} $ \\
\\[-0.3em]
$ \mu \geq 0.2 $ &
$ \forall i: \ s_i(t_\emptyset) = S_1 $ &
$ 2.5\mu - 1 $ \\
\bottomrule
\end{tabular}
\caption{Symmetric equilibrium strategies upon receiving $t_\emptyset$}
\label{tab:equilibrium-strategies}
\end{table}
\newline
Intuitively, when commitments occur simultaneously ($\mu = 0$), each player commits to going straight with probability $\tfrac{8}{11}$ and obtains an expected payoff of $-\tfrac{5}{11}$, matching the results from the Game of Chicken presented in the introduction. What happens as $\mu$ increases? Players gain more information about the order of commitments, which one might initially expect to reduce miscoordination and thus improve payoffs. 
However, this is not necessarily the case. For certain intermediate values of $\mu$, equilibrium payoffs become strictly lower than under simultaneous commitment, as illustrated in Figure~\ref{fig:mu-payoff}. Here, the dashed horizontal line represents the payoff under simultaneous commitment, while the solid line depicts equilibrium payoffs for the symmetric equilibrium across values of $\mu$. Notably, equilibrium payoffs dip below the simultaneous commitment payoff for values approximately in the interval $[0.11,0.22]$. 
Figure~\ref{fig:mu-miscoord-proba} further clarifies that in this interval, the probability of miscoordination is higher than under simultaneous commitment.
\footnote{Calculated as the sum of probabilities for miscoordinating directly on commitments and miscoordinating in the base game: $P(\omega_\emptyset)\cdot [p(S_1)^2 + (1 - p(S_1))^2 \cdot (\tfrac{2}{3})^2]$.}
Hence, more information about the order of commitment can lead to greater miscoordination and worse outcomes.
\begin{figure}[h!]
    %––––––––––– knobs you can play with ––––––––––––––––––––––
    \newcommand{\YMAXone}{-0.2}   % top of panel-A (orig. payoff)
    \newcommand{\XRMAXone}{0.3}   % right cut-off panel-A  (≥0.2)
    \newcommand{\YMINtwo}{0.4}    % bottom of panel-B (alt. payoff)
    \newcommand{\XRMAXtwo}{0.3}   % right cut-off panel-B  (≥0.2)
    %––––––––––––––––––––––––––––––––––––––––––––––––––––––––––
    
    \centering
    %--------------- PANEL A --------------------------------------
    \begin{minipage}[t]{0.48\linewidth}
        \centering
        \begin{tikzpicture}
            \begin{axis}[
                width=7.2cm,height=6cm,
                axis x line=bottom,
                axis y line=middle,
                axis line style={ultra thick},
                xlabel={$\mu$},
                ylabel={$U$},
                xmin=0, xmax=\XRMAXone,
                ymin=-0.55, ymax=\YMAXone,
                xtick={0,0.1,0.2,0.3,0.4,0.5}, % ticks > xmax are ignored
                ytick={-0.5,-0.3},
                tick label style={font=\small},
                label style={font=\small},
                clip=false]
            % mixed equilibrium (0 ≤ μ ≤ 0.2)
            \addplot[domain=0:0.2,samples=200,very thick,black,
                     restrict y to domain=-0.55:\YMAXone,
                     unbounded coords=discard]
              { (3*x - 5)/11 + (-2*x - 23*x*x)/(22*(1 - 2*x)) };
 
            % pure strategy     (0.2 ≤ μ ≤ XRMAXone)
            \addplot[domain=0.2:\XRMAXone,samples=200,very thick,black,
                     restrict y to domain=-0.55:\YMAXone,
                     unbounded coords=discard]
              { 2.5*x - 1 };
 
            % guides
            \addplot[dashed, thin] coordinates {(0.2,-0.55) (0.2,\YMAXone)};
            \node[font=\small, anchor=south]
                  at (axis cs:0.2,\YMAXone) {$\mu=0.2$};
 
            \addplot[dashed, thin] coordinates {(0,-5/11) (\XRMAXone,-5/11)};
            \node[font=\small, anchor=east]
                  at (axis cs:0,-5/11) {$-\tfrac{5}{11}$};
            \end{axis}
        \end{tikzpicture}
        \caption{Equilibrium payoff}
        \label{fig:mu-payoff}
    \end{minipage}
    \hfill
    %--------------- PANEL B --------------------------------------
    \begin{minipage}[t]{0.48\linewidth}
        \centering
        \begin{tikzpicture}
            \begin{axis}[
                width=7.2cm,height=6cm,
                axis x line=bottom,
                axis y line=middle,
                axis line style={ultra thick},
                xlabel={$\mu$},
                ylabel={$P$},
                xmin=0, xmax=\XRMAXtwo,
                ymin=\YMINtwo, ymax=0.65,
                xtick={0,0.1,0.2,0.3,0.4,0.5},
                ytick={\YMINtwo,0.5,0.6},
                tick label style={font=\small},
                label style={font=\small},
                clip=false]
            % mixed equilibrium (0 ≤ μ ≤ 0.2)
            \addplot[domain=0:0.2,samples=250,very thick,black,
                     restrict y to domain=\YMINtwo:0.65,
                     unbounded coords=discard]
              { (149*x*x - 152*x + 68) / (121*(1 - 2*x)) };
 
            % pure strategy     (0.2 ≤ μ ≤ XRMAXtwo)
            \addplot[domain=0.2:\XRMAXtwo,samples=200,very thick,black,
                     restrict y to domain=\YMINtwo:0.65,
                     unbounded coords=discard]
              { 1 - 2*x };
 
            % guides
            \addplot[dashed, thin] coordinates {(0.2,\YMINtwo) (0.2,0.65)};
            \node[font=\small, anchor=south]
                  at (axis cs:0.2,0.65) {$\mu=0.2$};
 
            \addplot[dashed, thin] coordinates {(0,68/121) (\XRMAXtwo,68/121)};
            \node[font=\small, anchor=east]
                  at (axis cs:0,68/121) {$\tfrac{68}{121}$};
            \end{axis}
        \end{tikzpicture}
        \caption{Miscoordination probability}
        \label{fig:mu-miscoord-proba}
    \end{minipage}
\end{figure}
\newline
To build intuition for this result, suppose player $i$ decides to wait. If $\mu = 0$, the state is $\omega_\emptyset$ and two outcomes are possible: either the other player simultaneously committed, causing $i$ to lose the Game of Chicken, or the other player also waited, leading to both playing the base game. When $\mu > 0$, a third outcome is possible: the state is $\omega_i$, meaning $i$ was actually the first mover. By waiting, $i$ allows the second player to observe her inaction and immediately respond by committing straight, causing $i$ to lose again. This makes waiting less attractive, prompting players to commit Straight more frequently. Consequently, the probability of miscoordination increases, lowering equilibrium payoffs. As $\mu$ grows larger, waiting eventually becomes dominated, leading to both committing straight and achieving the lowest possible payoff. However, as $\mu$ continues to increase, simultaneous commitments become rarer, and the probability of miscoordination decreases again.
 
This concludes the current section. I have employed a stylized Game of Chicken to illustrate the dynamics arising under uncertainty about the order of commitment. Extending this approach to a wider class of games is an interesting future direction but falls outside the scope of this thesis.
 

\end{document}